\documentclass[journal]{IEEEtran}
\usepackage[utf8]{inputenc}
\usepackage[T1]{fontenc}
\usepackage[spanish]{babel}
\addto\captionsspanish{\renewcommand{\abstractname}{Resumen}}
\usepackage{graphicx}
\usepackage{caption}
\usepackage{float}
\usepackage{url}
\usepackage{hyperref}

\hyphenation{op-tical net-works semi-}

\begin{document}

\title{Bit\'acora de Ghidra: CODEBREAKERS\\
\large El misterio del c\'odigo oculto: paso a paso del an\'alisis}

\author{Brayan Loaiza Leal,
Juan Sebasti\'an Guayaz\'an Clavijo
\thanks{Fundamentos de Seguridad de la Informaci\'on}
\thanks{Ingenier\'ia de Sistemas}
\thanks{Escuela Colombiana de Ingenier\'ia Julio Garavito}
\thanks{2026, Tercer Tercio}}

\markboth{Laboratorio 06 --- Bit\'acora de Ghidra}%
{Reto CODEBREAKERS}

\maketitle

\begin{abstract}
Este documento registra, en el orden en que se hizo, c\'omo se analiz\'o en Ghidra
el programa del reto CODEBREAKERS hasta encontrar la clave y obtener la FLAG, sin
acceso al c\'odigo fuente. Se parte de una ejecuci\'on con una clave falsa, se
importa el binario en Ghidra, se localizan las cadenas del programa y se sigue la
referencia al mensaje de acceso hasta \texttt{main}. All\'i se identifica un
descifrado AES-256-CBC cuya clave e IV est\'an en texto plano dentro del
ejecutable; se reproduce el descifrado con OpenSSL y se comprueba la clave
obtenida ejecutando el programa.
\end{abstract}

\begin{IEEEkeywords}
Ghidra, ingenier\'ia inversa, CTF, AES-256-CBC, an\'alisis est\'atico.
\end{IEEEkeywords}

\IEEEpeerreviewmaketitle

\section{Datos generales}

\begin{itemize}
  \item \textbf{Fecha:} 2026-11-11.
  \item \textbf{Herramienta:} Ghidra 12.1.4 (release oficial de la NSA) sobre
        JDK 26, en Windows 11. Proyecto \texttt{ECI\_\allowbreak{}CODEBREAKERS}.
  \item \textbf{Apoyo:} WSL2 Ubuntu 24.04 para ejecutar el programa y OpenSSL
        para reproducir el descifrado.
  \item \textbf{Binario:} \texttt{reto\_\allowbreak{}codebreakers\_\allowbreak{}linux\_\allowbreak{}x64}, ELF 64-bit PIE,
        x86-64, din\'amico, \textit{not stripped}.
  \item \textbf{SHA-256:}
\end{itemize}
{\scriptsize\begin{verbatim}
8676a8798847da67b0449092e4341eaaf21a995a37d7
aea916485b1503dc717f
\end{verbatim}}
\begin{itemize}
  \item \textbf{Resultado:} la clave \texttt{BREAKER\#2026} produce
        \texttt{ACCESO CONCEDIDO} y \texttt{FLAG\{CODEBREAKERS\_\allowbreak{}ECI\}}.
\end{itemize}

\subsection{Flujo general}

{\scriptsize\begin{verbatim}
Ejecutar (clave falsa)
 -> Importar en Ghidra -> Auto Analyze
 -> Search for Strings
 -> "ACCESO CONCEDIDO" -> References
 -> main -> Decompiler:
    strcmp(entrada, descifrado)
 -> que es "descifrado"?
    sbox / rsbox / add_round_key
    => AES-256-CBC
 -> clave e IV en 0x102300 y 0x102320,
    cifrado en aes_ct
 -> descifrar con OpenSSL
 -> BREAKER#2026 -> ejecutar -> FLAG
\end{verbatim}}

\section{Paso 0. Primer intento sin an\'alisis}

Antes de abrir Ghidra se ejecut\'o el programa en WSL con una palabra cualquiera
para conocer su comportamiento:

{\scriptsize\begin{verbatim}
$ chmod +x reto_codebreakers_linux_x64
$ ./reto_codebreakers_linux_x64
=== CODEBREAKERS: BOVEDA DIGITAL ECI ===
Desafio: descubre la clave sin recibir
el codigo fuente.
Introduce la clave: prueba
ACCESO DENEGADO
\end{verbatim}}

\noindent\textbf{Observaci\'on:} el programa pide la clave por la entrada est\'andar
(no como argumento) y responde \texttt{ACCESO DENEGADO}. El objetivo es encontrar
qu\'e texto hace que responda lo contrario.

\section{Paso 1. Crear el proyecto e importar}

\begin{table}[H]
\centering
\caption{Creaci\'on del proyecto e importaci\'on del binario.}
\footnotesize
\begin{tabular}{@{}l p{1.35in} p{1.45in}@{}}
\hline
\# & Acci\'on & Resultado \\ \hline
1.1 & Abrir \texttt{ghidraRun.bat} desde la carpeta de Ghidra & Ghidra 12.1.4 arranca con JDK 26. \\
1.2 & File $\rightarrow$ New Project $\rightarrow$ Non-Shared Project, nombre \texttt{ECI\_\allowbreak{}CODEBREAKERS} & Proyecto local, sin servidor. \\
1.3 & File $\rightarrow$ Import File\ldots{} $\rightarrow$ \texttt{reto\_\allowbreak{}codebreakers\_\allowbreak{}linux\_\allowbreak{}x64} & Formato ELF, lenguaje \texttt{x86:LE:64:default}, compilador gcc. Base de imagen \texttt{0x100000}. \\
1.4 & Doble clic en el programa $\rightarrow$ Analyze now? $\rightarrow$ Yes, opciones por defecto & An\'alisis autom\'atico de funciones, cadenas, referencias y decompilaci\'on. \\ \hline
\end{tabular}
\end{table}

\noindent\textbf{Observaci\'on:} el binario no est\'a \textit{stripped}. En
Symbol Tree $\rightarrow$ Functions aparecen nombres \'utiles: \texttt{main}
(\texttt{0x1010c0}) y \texttt{add\_\allowbreak{}round\_\allowbreak{}key} (\texttt{0x101ba0}), y en los datos
\texttt{sbox}, \texttt{rsbox} y \texttt{aes\_\allowbreak{}ct}.

\section{Paso 2. Buscar cadenas}

\begin{table}[H]
\centering
\caption{B\'usqueda de cadenas.}
\footnotesize
\begin{tabular}{@{}l p{1.1in} p{1.7in}@{}}
\hline
\# & Acci\'on & Resultado \\ \hline
2.1 & Search $\rightarrow$ For Strings\ldots{}, opciones por defecto & Lista de todas las cadenas del binario. \\
2.2 & Revisar los mensajes del programa & \texttt{ACCESO CONCEDIDO} (\texttt{0x1020a2}), \texttt{FLAG\{CODEBREAKERS\_\allowbreak{}ECI\}} (\texttt{0x1020b3}), \texttt{ACCESO DENEGADO} (\texttt{0x1020ca}). \\
2.3 & Buscar algo con forma de clave & Una cadena de 48 caracteres en \texttt{0x102300} (ver abajo). \\ \hline
\end{tabular}
\end{table}

{\scriptsize\begin{verbatim}
CodeBreakersECI_AES256_KEY_2026!
ECI_VECTOR_2026!
\end{verbatim}}

\noindent\textbf{Pista encontrada:} la cadena larga no es la clave del reto:
anuncia \texttt{AES256\_\allowbreak{}KEY} y \texttt{VECTOR}. Esto indica que la clave del reto
est\'a cifrada con AES-256 y que esa cadena contiene la clave AES y el vector de
inicializaci\'on (IV). Se verific\'o con \texttt{strings}:

{\scriptsize\begin{verbatim}
$ strings -n 6 reto_codebreakers_linux_x64 \
    | grep -E 'ACCESO|FLAG|CodeBreakers'
ACCESO CONCEDIDO
FLAG{CODEBREAKERS_ECI}
ACCESO DENEGADO
CodeBreakersECI_AES256_KEY_2026!ECI_VECTOR_2026!
\end{verbatim}}

\section{Paso 3. Seguir la referencia a ACCESO CONCEDIDO}

\begin{table}[H]
\centering
\caption{Seguimiento de la referencia hasta \texttt{main}.}
\footnotesize
\begin{tabular}{@{}l p{1.35in} p{1.45in}@{}}
\hline
\# & Acci\'on & Resultado \\ \hline
3.1 & Doble clic en \texttt{ACCESO CONCEDIDO} en la lista de cadenas & Se abre el Listing en \texttt{0x1020a2}. \\
3.2 & Clic derecho $\rightarrow$ References $\rightarrow$ Show References to Address & Una sola referencia, en \texttt{0x101a5b}, dentro de \texttt{main}. \\
3.3 & Doble clic en la referencia y abrir la ventana Decompiler & Se ve el \texttt{if} que decide entre los dos mensajes. \\ \hline
\end{tabular}
\end{table}

Final de \texttt{main} en el Decompiler:

{\scriptsize\begin{verbatim}
local_1c8 = _DAT_00102320 ^ local_1c8;
 // XOR con el IV (modo CBC)
uStack_1c0 = _UNK_00102328 ^ uStack_1c0;
bStack_1c9 = (byte)(uStack_1c0 >> 0x38);
 // ultimo byte = relleno
if ((byte)(bStack_1c9 - 1) < 0x10) {
  *(undefined1 *)((long)&local_1c8 +
     (long)(int)(0x10 - (uint)bStack_1c9)) = 0;
}
iVar21 = strcmp((char *)local_1a8,
                (char *)&local_1c8);
if (iVar21 == 0) {
  puts("ACCESO CONCEDIDO");
  puts("FLAG{CODEBREAKERS_ECI}");
}
else {
  puts("ACCESO DENEGADO");
}
\end{verbatim}}

\noindent\textbf{Conclusi\'on:} la entrada del usuario (\texttt{local\_\allowbreak{}1a8}, le\'ida
con \texttt{fgets}) se compara con \texttt{strcmp} contra un b\'ufer de 16 bytes
(\texttt{local\_\allowbreak{}1c8}). Ese b\'ufer no est\'a en claro en el binario: \texttt{main}
lo calcula antes de comparar. Hay que entender c\'omo.

\section{Paso 4. Identificar el algoritmo}

Se ley\'o \texttt{main} desde el inicio en el Decompiler, renombrando lo que se
iba entendiendo.

\begin{table}[H]
\centering
\caption{Lectura de \texttt{main} en el Decompiler.}
\footnotesize
\begin{tabular}{@{}l p{1.5in} p{1.3in}@{}}
\hline
\# & Qu\'e se observ\'o en Ghidra & Interpretaci\'on \\ \hline
4.1 & \texttt{local\_\allowbreak{}128 = \_\allowbreak{}DAT\_\allowbreak{}00102300} \ldots{} \texttt{uStack\_\allowbreak{}110 = \_\allowbreak{}UNK\_\allowbreak{}00102318}: se copian 32 bytes desde \texttt{0x102300} & Los primeros 32 caracteres de la cadena larga se cargan como clave. \\
4.2 & Bucle \texttt{do \ldots{} while (uVar23 != 0x3c)} que usa \texttt{sbox[...]} y \texttt{1 \textless{}\textless{} (\ldots)} cada 8 palabras & Expansi\'on de clave de AES-256: 60 palabras (\texttt{0x3c}) y la constante de ronda. \\
4.3 & \texttt{puVar19 = \&aes\_\allowbreak{}ct}; copia 16 bytes de \texttt{aes\_\allowbreak{}ct} (\texttt{0x1020e0}) & Es el texto cifrado: un solo bloque AES de 16 bytes. \\
4.4 & \texttt{puVar19 = rsbox;} y llamadas a \texttt{add\_\allowbreak{}round\_\allowbreak{}key(local\_\allowbreak{}1b8, \&local\_\allowbreak{}128, \ldots)} & \texttt{rsbox} es la S-box inversa de AES: el programa descifra, no cifra. \\
4.5 & \texttt{local\_\allowbreak{}1c8 = \_\allowbreak{}DAT\_\allowbreak{}00102320 \^{} local\_\allowbreak{}1c8} (16 bytes desde \texttt{0x102320}) & XOR con el IV tras descifrar el bloque: modo CBC. El IV es \texttt{ECI\_\allowbreak{}VECTOR\_\allowbreak{}2026!}. \\
4.6 & Se toma el \'ultimo byte y se pone un 0 en la posici\'on \texttt{0x10 - relleno} & Quita el relleno PKCS\#7 y termina la cadena para \texttt{strcmp}. \\ \hline
\end{tabular}
\end{table}

\begin{table}[H]
\centering
\caption{Datos identificados en el binario.}
\footnotesize
\begin{tabular}{@{}p{0.85in} l p{1.45in}@{}}
\hline
Dato & Direcci\'on & Valor \\ \hline
Clave AES-256 (32 B) & \texttt{0x102300} & \texttt{CodeBreakersECI\_\allowbreak{}AES256\_\allowbreak{}KEY\_\allowbreak{}2026!} \\
IV (16 B) & \texttt{0x102320} & \texttt{ECI\_\allowbreak{}VECTOR\_\allowbreak{}2026!} \\
\texttt{aes\_\allowbreak{}ct} (16 B) & \texttt{0x1020e0} & \texttt{d67f5789f451afeb} \texttt{f2b6e566f87fb979} \\
\texttt{rsbox} (256 B) & \texttt{0x102100} & S-box inversa de AES \\
\texttt{sbox} (256 B) & \texttt{0x102200} & S-box de AES, usada en la expansi\'on de clave \\ \hline
\end{tabular}
\end{table}

\section{Paso 5. Descifrar la clave}

Con la clave, el IV y el texto cifrado identificados en Ghidra, se reprodujo en
WSL el mismo descifrado que hace \texttt{main}:

{\scriptsize\begin{verbatim}
KEY=436f6465427265616b6572734543495f
    4145533235365f4b45595f3230323621
IV=4543495f564543544f525f3230323621
CT=d67f5789f451afebf2b6e566f87fb979
$ printf '%s' "$CT" | xxd -r -p | \
  openssl enc -d -aes-256-cbc -K $KEY -iv $IV
BREAKER#2026
\end{verbatim}}

OpenSSL retira el relleno PKCS\#7 igual que el programa. El resultado es la
clave: \texttt{BREAKER\#2026}.

\section{Paso 6. Comprobar ejecutando el programa}

{\scriptsize\begin{verbatim}
$ ./reto_codebreakers_linux_x64
=== CODEBREAKERS: BOVEDA DIGITAL ECI ===
Desafio: descubre la clave sin recibir
el codigo fuente.
Introduce la clave: BREAKER#2026
ACCESO CONCEDIDO
FLAG{CODEBREAKERS_ECI}
\end{verbatim}}

\begin{table}[H]
\centering
\caption{Comparaci\'on de las dos pruebas.}
\footnotesize
\begin{tabular}{@{}l l p{1.3in}@{}}
\hline
Prueba & Entrada & Salida \\ \hline
Clave falsa & \texttt{prueba} & \texttt{ACCESO DENEGADO} \\
Clave descubierta & \texttt{BREAKER\#2026} & \texttt{ACCESO CONCEDIDO} y \texttt{FLAG\{CODEBREAKERS\_\allowbreak{}ECI\}} \\ \hline
\end{tabular}
\end{table}

\section{Evidencia para el monitor}

\begin{table}[H]
\centering
\caption{Respuestas a la gu\'ia del reto.}
\footnotesize
\begin{tabular}{@{}p{1.1in} p{1.8in}@{}}
\hline
Pregunta de la gu\'ia & Respuesta \\ \hline
\textquestiondown Cu\'al es la clave? & \texttt{BREAKER\#2026} \\
\textquestiondown Qu\'e palabra o funci\'on llev\'o a la conclusi\'on? & La cadena \texttt{CodeBreakersECI\_\allowbreak{}AES256\_\allowbreak{}KEY\_\allowbreak{}2026!ECI\_\allowbreak{}VECTOR\_\allowbreak{}2026!}, la referencia de \texttt{ACCESO CONCEDIDO} hacia \texttt{main}, el \texttt{strcmp} de \texttt{main} y los s\'imbolos \texttt{rsbox}, \texttt{add\_\allowbreak{}round\_\allowbreak{}key} y \texttt{aes\_\allowbreak{}ct}, que identifican un descifrado AES-256-CBC. \\
\textquestiondown Qu\'e mensaje produce la clave correcta? & \texttt{ACCESO CONCEDIDO} seguido de \texttt{FLAG\{CODEBREAKERS\_\allowbreak{}ECI\}} \\ \hline
\end{tabular}
\end{table}

\section{Atajos de Ghidra usados}

\begin{table}[H]
\centering
\caption{Atajos y men\'us empleados.}
\footnotesize
\begin{tabular}{@{}p{1.5in} p{1.4in}@{}}
\hline
Atajo o men\'u & Acci\'on \\ \hline
Search $\rightarrow$ For Strings\ldots{} & Listar las cadenas del binario \\
Clic derecho $\rightarrow$ References $\rightarrow$ Show References to Address & Ver desde d\'onde se usa una cadena o un dato \\
Doble clic & Seguir una referencia, una llamada o un dato \\
\texttt{L} & Renombrar variable, funci\'on o etiqueta \\
\texttt{;} & Agregar comentario \\
\texttt{G} & Ir a una direcci\'on, por ejemplo \texttt{0x102300} \\ \hline
\end{tabular}
\end{table}

\section{Reflexi\'on}

\begin{enumerate}
  \item \textbf{La clave viaja con el programa.} La clave AES y el IV est\'an
        dentro del mismo ejecutable, en texto plano. Cifrar la contrase\~na solo
        a\~nade un paso: cualquiera que lea el binario puede descifrarla.
  \item \textbf{La FLAG no necesita la clave.}
        \texttt{FLAG\{CODEBREAKERS\_\allowbreak{}ECI\}} est\'a en texto plano y aparece con un
        simple \texttt{strings}.
  \item \textbf{Los s\'imbolos ayudan al analista.} Al no estar \textit{stripped},
        nombres como \texttt{rsbox}, \texttt{add\_\allowbreak{}round\_\allowbreak{}key} y \texttt{aes\_\allowbreak{}ct}
        delatan el algoritmo de inmediato.
  \item \textbf{Ofuscar no es proteger.} Un secreto que el programa necesita para
        decidir debe quedarse en un servidor, no en el cliente.
\end{enumerate}

\end{document}
